\documentclass[11pt]{amsart}
\usepackage[margin=1.1in]{geometry}
\usepackage{amsmath,amssymb,amsthm}
\usepackage{booktabs}
\usepackage[dvipsnames]{xcolor}
\usepackage{hyperref}
\hypersetup{colorlinks=true,linkcolor=blue!50!black,citecolor=blue!50!black,urlcolor=blue!50!black}

\newtheorem{theorem}{Theorem}
\newtheorem{proposition}[theorem]{Proposition}
\newtheorem{lemma}[theorem]{Lemma}
\theoremstyle{definition}
\newtheorem{definition}[theorem]{Definition}
\theoremstyle{remark}
\newtheorem{remark}[theorem]{Remark}

\newcommand{\Z}{\mathbb{Z}}
\newcommand{\up}{{\uparrow}}

\title{A covering system with minimum modulus 44}
% TODO(author): confirm name, affiliation and e-mail.
\author{Grayson MacKenzie}
\date{September 2026}

\def\UrlBreaks{\do\/\do-\do.}
\begin{document}

\begin{abstract}
We construct a covering system of the integers with pairwise distinct moduli whose least modulus is $44$. Nielsen (2009) constructed such a system with least modulus $40$, and Owens (2014) extended the method to $42$.
Starting from a variant of Owens's construction, we extend it twice, first to least modulus $43$ and then to $44$. The new material
consists of roots on fresh primes up to $211$, small ``device'' families, and in-place repairs found by an exact
counterexample loop. The final system consists of $618{,}413$ families of congruences on $47$ primes; after a
standard finite truncation it is a finite covering system. Coverage was decided exactly, twice, by independent SAT
encodings whose unsatisfiability proofs were checked by independent proof checkers.
\end{abstract}

\maketitle

\section{Introduction}

A \emph{covering system} is a finite set of congruence classes $a_1 \bmod m_1, \dots, a_k \bmod m_k$ whose union is
$\Z$. Erdős introduced covering systems with distinct moduli in 1950 \cite{Erdos50}. He asked whether the least
modulus of such a system can be arbitrarily large, a question he later described as perhaps his favourite problem.
Hough \cite{Hough15} answered it in the negative, building on Filaseta, Ford, Konyagin, Pomerance and Yu
\cite{FFKPY07}: every covering system with distinct moduli has least modulus at most $10^{16}$. Balister, Bollobás,
Morris, Sahasrabudhe and Tiba \cite{BBMST22} gave a simpler proof and lowered the bound to $616{,}000$.

Constructions lag far behind this bound. Successive records raised the least modulus to $25$ (Gibson
\cite{Gibson09}) and to $40$ (Nielsen \cite{Nielsen09}); see \cite{Nielsen09} for the history. Owens
\cite{Owens14} extended Nielsen's method to $42$ in a master's thesis.

\begin{theorem}\label{thm:main}
There is a covering system of $\Z$ with pairwise distinct moduli whose least modulus is $44$.
\end{theorem}

\noindent Our contributions are as follows.
\begin{enumerate}
\item \emph{The starting system} (§\ref{sec:owens}). We work with UC85, a variant of Owens's construction that
  differs from it in the placement of two carrier congruences and in the arrangement of four later sections. It
  has $356{,}364$ families and least modulus $42$.
\item \emph{From $42$ to $43$ to $44$} (§§\ref{sec:43}–\ref{sec:44}). We cut the unique modulus-$42$ family back to
  its deeper levels. We then re-cover exactly what it alone covered, using roots on fresh primes $89$–$173$, a small
  device and in-place fills. We do the same for the unique modulus-$43$ family, with a $41$-arrow in its slot, slot
  fills, roots on fresh primes $179$–$211$ and $1{,}801$ repairs found by an exact counterexample loop. Every modulus
  from $44$ to $160$ except $121$ is used exactly once in the final system.
\item \emph{Verification} (§\ref{sec:verify}). Two independent exact checks establish coverage. Each builds a finite
  propositional formula that is satisfiable exactly when some integer other than a ray centre is uncovered (the
  single ray centre left is covered by the finite completion of §\ref{sec:prelim}); the formula is unsatisfiable,
  and every step of the solver's refutation is checked by separate programs. Both checks also confirm distinct
  moduli.
\end{enumerate}

\noindent The complete list of families and the certificates are available at
\url{https://github.com/gr-mk/covering-minimum-modulus-44}.

\section{Families, distinct moduli and finite completion}\label{sec:prelim}

\subsection*{Selectors and families}
We describe congruence classes through \emph{selectors} on single primes. For a prime $p$:
\begin{itemize}
\item a \emph{fixed} selector $p{:}\,r \bmod p^k$ matches $x$ when $x \equiv r \pmod{p^k}$;
\item a \emph{ray} selector $\mathrm{ray}(p,[\ell,h],c,d)$, with $1\le \ell\le h\le\infty$, centre $c\in\Z$ and
  digit $d\in\{1,\dots,p-1\}$, matches $x \ne c$ when $e := 1+v_p(x-c)$ lies in $[\ell,h]$ and
  $(x-c)/p^{e-1} \equiv d \pmod p$.
\end{itemize}
Equivalently, the ray matches the classes $x \equiv c + d\,p^{e-1} \pmod{p^e}$ for $\ell\le e\le h$. A \emph{family}
is a product of selectors on distinct primes, and it matches $x$ when every selector does. Choosing one level $e_p$
for each ray, a family is a union of single residue classes, one for each exponent vector, with modulus
$\prod_p p^{e_p}$ (where $e_p=k$ for a fixed selector). These exponent vectors form the family's \emph{exponent box}.

\subsection*{Distinct moduli}
Every modulus of a family is divisible by exactly the primes of its support. Two families with different supports
therefore never share a modulus, and the moduli of a system are pairwise distinct exactly when, among families with
a common support, the exponent boxes are pairwise disjoint.

\subsection*{Ray centres and finite completion}
A ray never matches its own centre. Consequently a system of families can cover every integer except finitely many
\emph{ray centres}. Our final system misses only the integer $0$.

Rays are infinite recursions, and we make the system finite as Nielsen does \cite[§2]{Nielsen09}. Each ray is
truncated at a sufficiently deep level, and the finitely many deep classes that remain, including the ray centre
$0$, are covered by congruences that also involve an auxiliary prime $P$ dividing no other modulus. Distinct powers
of the prime of the ray (and of $P$ where necessary) keep all moduli distinct, and every new modulus is a multiple
of $P$, so it exceeds $44$. We take $P=223$, the least prime not used by the system.

\section{The starting system}\label{sec:owens}

We use the notation of Nielsen \cite{Nielsen09} and Owens \cite{Owens14}. A $q$-arrow
$q\up(A_1,\dots,A_{q-1})$ covers a region by assigning to each nonzero $q$-adic digit $j$ an \emph{input} $A_j$,
itself a set of congruences (with multipliers) that covers the region inside that digit's slot. Inputs of the form
$m$ (a bare multiplier) give single congruences whose modulus is $q$ times $m$ at the first level.

Owens's construction has the same shape as Nielsen's:
\begin{itemize}
\item a base on the primes $2,3,5$ (§§3.1–3.3);
\item sections for the primes $7$ to $83$ (§§3.4–3.20, the last three holding two primes each), whose arrows
  cover the parts of $\Z$ the earlier sections leave open.
\end{itemize}
The least modulus $42 = 2\cdot3\cdot7$ is attained once, by the family $2{:}0 \bmod 2 \cdot 3{:}1 \bmod 3 \cdot
\mathrm{ray}(7,[1,\infty),0,3)$ of §3.4.

\subsection*{The variant UC85}
UC85 differs from Owens's construction \cite{Owens14} in three ways.
\begin{enumerate}
\item The modulus-$55$ congruence of §3.5, the bare input $1$ of the sixth input $5\up(1,2,3\cdot1,4)$ of the
  $11$-arrow, is placed at $1 \bmod 5$, as in Nielsen's construction.
\item The modulus-$85$ congruence of §3.7, the bare $1$ in a $5\up(1,2,4,8\up)$ input of its $17$-arrow, is placed at
  $4 \bmod 5$ by permuting that $5$-arrow to $5\up(8\up,2,4,1)$. This gives §3.18's $67$-arrow a carried
  $17$-slot, and §3.18 then needs no $89$-arrow.
\item The arrow layers of §§3.12, 3.14, 3.15 and 3.18 are arranged by an exact planner (a subset dynamic
  programme over the available input sets).
\end{enumerate}
The result has $356{,}364$ families on the $23$ primes $2,\dots,83$. Its least modulus is $42$, attained once. It was
certified by the method of §\ref{sec:verify}.

\section{From 42 to 43}\label{sec:43}

Let $F_{42}$ be the modulus-$42$ family above. Changing its ray to $\mathrm{ray}(7,[2,\infty),0,3)$ removes exactly the
level-one class $10 \bmod 42$ and keeps the moduli $2\cdot3\cdot7^e$ for $e\ge2$. Most of $10 \bmod 42$ is still
covered by other families. The exact uncovered part $U$ is $2.1\%$ of the class, about $4.95\cdot10^{-4}$ of $\Z$,
spread over $14$ cells of the form ($2$-adic valuation, residue mod~$5$). We re-cover $U$ with three ingredients:

\begin{itemize}
\item \textbf{A $7$-device.} Seven families on $\mathrm{ray}(7,[2,\infty),-4,j)$ inside $3 \bmod 7$, with moduli
  from $49$ to $23{,}520$. They cover fixed $7$-digits of $x+4$ exactly, and so act as exact carries for deeper
  covers.
\item \textbf{In-place fills.} $604$ families placed in empty slots of §§3.10, 3.12 and 3.15 (respectively $3$,
  $75$ and $526$).
\item \textbf{Seventeen roots on fresh primes} $89$, $97$, $101$, $103$, $107$, $109$, $113$, $127$, $131$, $137$, $139$,
  $149$, $151$, $157$, $163$, $167$ and $173$. Each root $q\up$ covers one pinned piece of $U$ with $q-1$ inputs. The inputs are built from the
  other primes together with carries certified exactly: the device's digits, and root digits proved covered by an
  exact gate.
\end{itemize}

The resulting system has $540{,}826$ families on $40$ primes. Its least modulus is $43$, attained once by the
digit-$1$ input of §3.14's $43$-arrow. The moduli $43,\dots,48$ each occur once. Only $F_{42}$ was changed and
nothing was removed; this was checked independently.

\section{From 43 to 44}\label{sec:44}

The unique modulus-$43$ family is the multiplier-$1$ input $\mathrm{ray}(43,[1,\infty),0,1)$ of §3.14. We replace it
by its tail $\mathrm{ray}(43,[2,\infty),0,1)$, whose least modulus is $1{,}849$, and re-cover what its first level
alone covered.
\begin{enumerate}
\item \textbf{A $41$-arrow in its slot.} In §3.14 the slot is filled by a $41$-arrow built from the section's other
  inputs ($177$ families). §3.14's region then closes exactly.
\item \textbf{Slot fills.} The removed family also carried the $43$-slot $1$ of every inner $43$-arrow elsewhere. We
  fill slot $1$ of all $264$ affected $43$-arrow copies in §§3.15–3.19 ($1{,}183$ families) and add $2$ families at
  one input of §3.20's $83$-arrow.
\item \textbf{Exact census.} What is still uncovered lies in $1{,}282$ pieces, all inside $22$ root inputs of the
  $67$- and $73$-arrows of §§3.18–3.19. Their total measure is estimated, by sampling, at about $10^{-8}$ of $\Z$. The census is certified
  complete by a gate formula that is unsatisfiable with all pieces blocked ($160$ sub-cubes, all with checked
  proofs).
\item \textbf{Closing it.} Seven roots on the fresh primes $179, 181, 191, 193, 197, 199, 211$ ($74{,}424$ families)
  cover the heavy inputs. $1{,}801$ further families come from an exact counterexample loop: each uncovered point
  is traced to the input that misses it, and that input receives one thin family or a small local arrow, checked for
  collisions against the whole system.
\end{enumerate}

The final system has $618{,}413$ families on the $47$ primes $2,\dots,211$. Its least modulus is $44$, attained
once by $2{:}1 \bmod 4 \cdot \mathrm{ray}(11,[1,\infty),0,1)$ of §3.5. Table~\ref{tab:ladder} shows how densely the
smallest moduli are packed.

\begin{table}[ht]
\centering
\caption{Moduli $m\le160$ of the final system. Each is used by exactly one family. The only integer in $[44,160]$
that is not a modulus is $121$.}\label{tab:ladder}
\begin{tabular}{@{}ll@{}}
\toprule
range & moduli used\\
\midrule
$44$–$120$ & all $77$ integers\\
$121$ & not used\\
$122$–$160$ & all $39$ integers\\
\bottomrule
\end{tabular}
\end{table}

\section{Verification}\label{sec:verify}

\subsection*{Exact coverage}
For each prime $p$ of the system, the selectors on $p$ partition $\Z$ into finitely many \emph{local atoms}. These
are residue classes modulo a fixed power of $p$, plus, for each ray, one atom per level up to a cap and one
``deep'' atom for the levels beyond it; a ray centre is its own limit atom. An integer's atoms on the different
primes determine which families contain it. Conversely, by the Chinese remainder theorem, every combination of atoms
is realised by some integer. Coverage of $\Z$ outside the ray centres is therefore equivalent to the
unsatisfiability of a propositional formula. That formula has one variable per atom, one exactly-one constraint per
prime, and one clause per family forbidding its atoms.

For the final system the formula has $618{,}804$ clauses. CaDiCaL \cite{CaDiCaL} proves it unsatisfiable, and two
further steps check that refutation:
\begin{itemize}
\item drat-trim \cite{DRATtrim} verifies the DRAT proof;
\item the trimmed LRAT proof is replayed step by step as pure reverse unit propagation ($77.7$ million hints).
\end{itemize}

A second check shares no code with the first:
\begin{itemize}
\item It reads the system with inert stand-ins, so none of the builder's code runs.
\item It uses two separately written encoders. Their refutations were checked by drat-trim, by lrat-check, and by
  kissat \cite{Kissat} combined with drat-trim.
\item Mutation controls confirm that neither encoder is vacuous. Removing a single needed root, repair or fill family
  makes the formula satisfiable, and the decoded model is an integer that only the removed family covers.
\end{itemize}

\subsection*{Distinct moduli and least modulus}
The second check confirms distinct moduli in two ways:
\begin{itemize}
\item a sweep over exponent boxes within each support;
\item a brute-force expansion of every modulus up to $10^{12}$, with no repeats and none below $44$.
\end{itemize}
Every family is well formed: its primes are sorted with one selector each, residues and digits are in range, and
$\ell\ge1$. A total of $50{,}912{,}769$ sampled integers (not a proof) were all covered, except the ray centre $0$.

\section{Remarks on 45 and beyond}\label{sec:beyond}

The same method does not obviously reach $45$. That would require removing the unique modulus-$44$ family, of §3.5.
\begin{itemize}
\item \emph{Patching.} Truncating that family opens about $3\cdot10^{-4}$ of $\Z$ in $78$ separate cells. A root on a
  fresh prime can cover only about one cell, and only primes up to about $257$ have enough room.
\item \emph{Re-planning inside the original regions.} Exact capacity certificates rule out re-planning §3.5 inside its
  region with the primes the system already uses. The shortfall is $5.3\%$ at the first $11$-level and $2.5\%$ for
  §§3.5, 3.7 and 3.11 jointly.
\item \emph{The remaining route.} Redrawing the region boundaries between §§3.5–3.7 is not excluded. With a new kind
  of root on free prime-pair copies of the lattice, everything except the region of a re-planned §3.6 can be covered
  exactly; the remaining integer problem is open.
\end{itemize}
Beyond $45$, the base on $2,3,5$ is itself saturated: it uses every $\{2,3,5\}$-smooth modulus at least $43$ exactly
once. Any construction of this type reaching $46$ or more needs a different base.

\section*{Acknowledgments}
The search and the certification described here were carried out with extensive assistance from
Claude, an AI model developed by Anthropic, working under the author's direction on the author's computing
resources. All results were checked by the exact methods of §\ref{sec:verify}.

\begin{thebibliography}{99}

\bibitem{BBMST22} P.~Balister, B.~Bollobás, R.~Morris, J.~Sahasrabudhe and M.~Tiba, \emph{On the Erdős covering
problem: the density of the uncovered set}, Invent. Math. \textbf{228} (2022), 377--414.

\bibitem{CaDiCaL} A.~Biere, K.~Fazekas, M.~Fleury and M.~Heisinger, \emph{CaDiCaL, Kissat, Paracooba, Plingeling and
Treengeling entering the SAT Competition 2020}, Proc. SAT Competition 2020.

\bibitem{Kissat} A.~Biere and M.~Fleury, \emph{Gimsatul, IsaSAT and Kissat entering the SAT Competition 2022}, Proc.
SAT Competition 2022.

\bibitem{Erdos50} P.~Erdős, \emph{On integers of the form $2^k+p$ and some related problems}, Summa Brasil. Math.
\textbf{2} (1950), 113--123.

\bibitem{FFKPY07} M.~Filaseta, K.~Ford, S.~Konyagin, C.~Pomerance and G.~Yu, \emph{Sieving by large integers and
covering systems of congruences}, J. Amer. Math. Soc. \textbf{20} (2007), 495--517.

\bibitem{Gibson09} D.~J.~Gibson, \emph{A covering system with least modulus 25}, Math. Comp. \textbf{78} (2009),
1127--1146.

\bibitem{Hough15} R.~Hough, \emph{Solution of the minimum modulus problem for covering systems}, Ann. of Math. (2)
\textbf{181} (2015), 361--382.

\bibitem{Nielsen09} P.~P.~Nielsen, \emph{A covering system whose smallest modulus is 40}, J. Number Theory
\textbf{129} (2009), 640--666.

\bibitem{Owens14} T.~E.~Owens, \emph{A covering system with minimum modulus 42}, Master's thesis, Brigham Young
University, 2014.

\bibitem{DRATtrim} N.~Wetzler, M.~J.~H.~Heule and W.~A.~Hunt~Jr., \emph{DRAT-trim: efficient checking and trimming
using expressive clausal proofs}, in: SAT 2014, Lecture Notes in Comput. Sci. \textbf{8561}, Springer, 2014,
422--429.

\end{thebibliography}

\end{document}
